% ═══════════════════════════════════════════════════════════════════════════
% CAPÍTULO 2 — REVISÃO DA LITERATURA (EXPANDIDO +4000 palavras)
% ═══════════════════════════════════════════════════════════════════════════
\chapter{Revisão da Literatura}
\label{ch:literatura}

\section{Metacognição: Fundamentos Teóricos}

\subsection{O Que É Metacognição?}

Metacognição, literalmente ``cognição sobre a cognição'', é a capacidade de refletir sobre os próprios processos mentais. Flavell (1979), pioneiro no estudo da metacognição, a definiu como ``o conhecimento e a regulação dos próprios processos cognitivos''. Para um leitor não familiarizado com o termo: assim como você pode estar lendo este texto (cognição) e simultaneamente perceber que está entendendo ou não o que lê (metacognição), um sistema de software pode executar uma tarefa e simultaneamente monitorar se está executando corretamente.

Flavell distinguiu dois componentes: \textbf{conhecimento metacognitivo} (o que sabemos sobre nosso funcionamento cognitivo — ``sou melhor em matemática que em poesia'') e \textbf{regulação metacognitiva} (como monitoramos e controlamos nossos processos — ``devo fazer um diagrama para entender melhor''). Esta distinção é fundamental porque mapeia diretamente para a arquitetura do OpenCode: o \texttt{SelfModel} implementa conhecimento metacognitivo (auto-representação), enquanto o \texttt{MetacognitiveMonitor} implementa regulação metacognitiva (auto-observação e correção).

\subsection{Metacognição em Inteligência Artificial}

Cox (2005) foi pioneiro na adaptação do conceito para IA. Em ``Metacognition in Computation'', argumentou que sistemas metacognitivos deveriam ser capazes de: (1) auto-monitoramento — observar desempenho e detectar desvios; (2) auto-avaliação — julgar qualidade dos outputs sem oráculo externo; (3) auto-regulação — ajustar processos internos baseado na auto-avaliação. O OpenCode implementa estes três níveis: o \texttt{AnomalyDetector} (monitoramento), o \texttt{ConfidenceEstimator} (avaliação), e o \texttt{CorrectionEngine} (regulação).

Cox e Raja (2011) posteriormente propuseram uma taxonomia de arquiteturas metacognitivas em três categorias: introspectivas (mantêm modelo explícito do próprio funcionamento), reflexivas (modificam seus próprios processos), e generativas (criam novas capacidades). O OpenCode se enquadra primariamente na categoria reflexiva (N2-N3), com potencial para evoluir para generativa (SPEC-034).

\subsection{A Taxonomia de Minsky}

Minsky (2006), em ``The Emotion Machine'', propôs uma arquitetura da mente com múltiplos níveis de reflexão. Cada nível observa e critica o inferior: Nível 0 (reativo: respostas instintivas), Nível 1 (aprendido: comportamentos por experiência), Nível 2 (deliberativo: planejamento consciente), Nível 3 (reflexivo: reflexão sobre o próprio pensamento), Nível 4 (auto-consciente: reflexão sobre a reflexão), Nível 5 (auto-idealização: valores e ética). Esta visão hierárquica inspirou diretamente a taxonomia N0-N3.5 implementada no OpenCode (Capítulo 6).

\section{Teorias de Consciência e Adaptação Computacional}

\subsection{Global Workspace Theory (GWT)}

A GWT, proposta por Baars (1988, 1997), postula que a consciência emerge de um ``espaço de trabalho global'' onde informações de múltiplos processadores especializados competem por acesso. A informação que ``vence'' é broadcastada para todo o sistema. O OpenCode implementa uma adaptação computacional: múltiplos módulos competem por acesso ao \texttt{GlobalWorkspace}; o \texttt{AttentionBuffer} (7 itens, Lei de Miller) seleciona qual informação entra; a informação selecionada é broadcastada para módulos registrados.

A evidência empírica para a GWT vem de estudos de neuroimagem (Dehaene \& Changeux, 2011) mostrando que estímulos conscientes ativam uma rede distribuída (o ``workspace global''), enquanto estímulos subliminares ativam apenas áreas locais. A adaptação computacional do OpenCode captura este princípio: informações no \texttt{GlobalWorkspace} são acessíveis a todos os módulos (``conscientes''), enquanto informações locais não são.

\subsection{Attention Schema Theory (AST)}

Graziano (2013, 2019) propôs que a consciência não é o processamento de informação em si, mas um \textbf{modelo} (schema) que o cérebro constrói para representar sua própria atenção. O OpenCode implementa a AST via \texttt{SelfModel}: mantém modelo de estados internos (confiança, anomalias, goals), forecasting preditivo (regressão linear sobre 10 snapshots), distinção self/other (classifica eventos como internos/externos), e histórico de 50 snapshots para análise de tendências.

\subsection{Atkinson-Shiffrin Memory Model}

Atkinson \& Shiffrin (1968) propuseram três estágios de memória: sensorial (TTL curto, alta capacidade), curto prazo (TTL médio, capacidade 7$\pm$2), e longo prazo (TTL longo, capacidade ilimitada). O OpenCode implementa via \texttt{NaturalForgetting} (SPEC-038): sensory (30s, 50 itens), short\_term (5min, 7 itens), long\_term (24h+, promoção por 5+ acessos com alta importância).

\section{Teoria dos Jogos e Governança de Commons}

\subsection{Ostrom's Design Principles (DP1-DP8)}

Elinor Ostrom (1990, 2010), Nobel de Economia 2009, desafiou a ``tragédia dos commons'' de Hardin (1968) demonstrando que comunidades autogerem recursos comuns por séculos através de 8 princípios de design: limites claros (DP1), regras proporcionais (DP2), participação coletiva (DP3), monitoramento (DP4), sanções graduais (DP5), resolução de conflitos (DP6), autonomia reconhecida (DP7), e empreendimentos aninhados (DP8). O OpenCode adapta estes princípios para governança de goals autônomos (Seção 6.3).

A evidência empírica para os princípios de Ostrom é robusta: Cox, Arnold \& Villamayor-Tomás (2010) realizaram uma meta-análise de 91 estudos e encontraram que os princípios DP1, DP2 e DP4 eram particularmente fortes preditores de sucesso na governança de commons. Isto informou a decisão de design de dar peso maior a estes três princípios na auditoria de goals do OpenCode.

\subsection{Equilíbrio de Nash e Estratégias Evolutivamente Estáveis}

O equilíbrio de Nash (Nash, 1950) — ponto onde nenhum agente tem incentivo para mudar unilateralmente — é aplicado no \texttt{DialecticalEngine} como ponto de convergência entre tese e antítese. Estratégias Evolutivamente Estáveis (Smith \& Price, 1973) — resistentes à invasão por mutantes — são aplicadas na seleção de correções: o sistema prefere estratégias com alta taxa de sucesso.

\section{Engenharia de Software com Agentes}

\subsection{Test-Driven Development como Método Científico}

Beck (2003) formalizou o TDD como ciclo Red-Green-Refactor. Janzen \& Saiedian (2005) demonstraram que TDD reduz defeitos em 40-90\% comparado com desenvolvimento tradicional. O OpenCode estende esta interpretação: cada CT é uma hipótese falseável (Popper, 1959); 312 CTs representam 312 hipóteses corroboradas, resistentes a 23 ciclos de refutação.

\subsection{Padrões de Projeto}

Gamma et al. (1994) catalogaram 23 padrões. O OpenCode utiliza 8 (Observer, Strategy, Chain of Responsibility, Factory, Singleton, Adapter, Decorator, Composite) e introduz 4 originais: Behavioral Gate (SPEC-038), Structural Compression (SPEC-037), Dialectical Synthesis (SPEC-036), e Ostrom Governance (SPEC-036).

\section{Trabalhos Relacionados}

\subsection{NEXO Brain}

O NEXO Brain (wazionapps/nexo, 2026) implementa ``cérebro compartilhado'' com metacognitive guard, trust scoring e natural forgetting. Três padrões foram adaptados: trust scoring com blend 70/30, behavioral gate com classificação de risco, e modelo Atkinson-Shiffrin.

\subsection{Awesome Autoresearch}

O Awesome Autoresearch (alvinreal/awesome-autoresearch, 2026) cataloga sistemas de auto-melhoria inspirados em Karpathy. O padrão ``experimento $\rightarrow$ avaliação $\rightarrow$ aprendizado'' inspirou o \texttt{OutcomeTracker} (SPEC-038).

\subsection{Ruflo e Planning-with-Files}

Ruflo (ruvnet/ruflo, 2026, 58.9k stars) oferece swarm intelligence com self-learning. Planning-with-files (OthmanAdi/planning-with-files, 2026, 23k stars) introduz planejamento determinístico baseado em arquivos para agentes. Ambos influenciaram decisões de design sobre persistência e coordenação multi-agente no OpenCode.
